\subsection{How the closure test drove the M2 iterations}
\label{sec:m2-iter}

The closure classifier turned out to be the most useful diagnostic of the project so far: each of its
failures pointed at a specific defect in how the reco variables were presented to the flow, not at the
network. Table~\ref{tab:m2-iter} lists the iterations, all trained on the same 32 files with the same encoder.

\begin{table}[h]\centering\small
\resizebox{\textwidth}{!}{\begin{tabular}{lcccccc}
\toprule
Iteration & AUC (reco only) & AUC (truth+reco) & $W_1$ $\log P$ ratio & $W_1$ $\Delta\theta_x$ & $W_1$ $\Delta$vtx$_z$ & $W_1$ $\log E_\mathrm{recoil}$ \\
\midrule
v1: angle differences, robust standardisation, clip $\pm 12\sigma$, 11 variables & 0.862 & 0.887 & 1.581 & 1.099 & 0.413 & 0.048 \\
v2: + asinh tail compression & 0.766 & 0.845 & 0.326 & 0.382 & 0.155 & 0.085 \\
v3: + calorimetric variants as log-ratios & 0.970 & 0.974 & 0.347 & 0.413 & 0.163 & 0.101 \\
v4: 9 generated + 2 derived, larger flow & \multicolumn{6}{c}{see Tables~\ref{tab:m2-tier1}, \ref{tab:m2-closure}} \\
\bottomrule
\end{tabular}
}
\caption{M2 iterations. $W_1$ values are in the model space of each iteration and are therefore only roughly
comparable across rows; the classifier AUCs are directly comparable.}
\label{tab:m2-iter}
\end{table}

\paragraph{v1: clipping created spikes.} The first version standardised each variable by a robust width and
clipped at $\pm 12\sigma$. The robust width of the muon variables is set by the narrow MINOS-matched core,
so the broad unmatched-muon population landed on the clip boundaries: 9\% of momentum ratios, 17\% of angle
differences and 5--9\% of vertex residuals were clipped, and among unmatched muons 76\% of the angle
differences. The flow was asked to reproduce artificial delta functions at the edges and could not.
The same run also exposed eight reconstructed events (of 4.8 million) whose tuple entries contain NaN or
infinite muon momenta and vertices; one such event per batch made the loss NaN, which silently prevented any
checkpoint from being saved in the very first attempt. These rows are now flagged and excluded from the flow
loss.

\paragraph{v2: asinh compression.} Replacing the clip by $\mathrm{asinh}(u/s_0)$ on the muon and vertex
columns made the tails logarithmic, removed all clipping and brought the AUC from 0.86 to 0.77. Per-variable
classifiers then showed every single marginal at AUC $\le 0.56$, but the calorimetry block alone at 0.80:
the discrimination came from \emph{pairs} of variables, specifically \texttt{recoil\_passivecorrected} versus
\texttt{recoil\_E} and \texttt{hadron\_recoil} versus \texttt{recoil\_E}.

\paragraph{v3: log-ratios, and a point mass.} Modelling those two as log-ratios to \texttt{recoil\_E} made
the problem worse (AUC 0.97, the passive-correction ratio alone 0.89), which revealed the actual structure:
the ratios are not smooth. The passive-corrected recoil equals $0.6669 \times$\texttt{recoil\_E} exactly for
more than 15\% of events (its 84th and 99th percentiles coincide) and takes a geometry-dependent value between
0.62 and 0.67 otherwise; \texttt{hadron\_recoil} equals $1.385 \times$\texttt{recoil\_E} to within a MeV for a
large fraction of events. These are deterministic MasterAnaDev calibrations of one number, not independent
detector responses, and a continuous flow cannot produce a delta function.

\paragraph{v4: derive, do not generate.} The two calibrated variants are removed from the generative target
and filled at decode time as \texttt{recoil\_E} times the median ratio in 250\,mm bins of true vertex $z$,
fitted on the training split. The approximation drops the event-by-event scatter of the calibration
(16--84\% widths of about 3\% and 10\% respectively) and is documented as such; an analysis that depends on
that scatter should use \texttt{recoil\_E} directly. The flow now generates 9 variables and was given more
capacity (5 layers of 768 units).
